% Include with % Include with % Include with % Include with \input{abstract} before Chapter 1 in the parent thesis.
% Numerical claims follow the frozen Chapter 6 60 Hz v5 formal results.
% This fragment has no dependency on chapter-specific result macros.

\chapter*{Abstract}

Live music-driven humanoid dance requires musical relevance, responsive timing,
continuous motion transitions, and feasible robot configurations to be maintained
under causal audio access and limited computation. This dissertation presents
Humanoid Matcher, a retrieval-and-control system that adapts a finite library of
dance trajectories to streaming music on a 29-degree-of-freedom Unitree G1 model.

The offline pipeline combines paired AIST++ music and human motion, extracts
salient key poses from whole-body velocity minima, and retargets SMPL sequences
through General Motion Retargeting. Provenance checks and MuJoCo kinematic
preflight produce a catalogue of 411 admitted motions associated with 60 music
identities and 348 audio segments. Online, six-second rolling audio windows are
gain-normalised and encoded using Discogs EffNet embeddings, style activations,
and rhythm--timbre descriptors. Hierarchical retrieval combines musical similarity
with tempo, key-pose-density, and activity compatibility. Temporal evidence and
bounded-hold selection govern motion changes, while asynchronous retrieval and
candidate preparation support continuous playback. Motion-aware beat acceptance
and bounded phase correction adapt timing gradually. Transition-entry search
uses joint pose, velocity, foot contact, and root state, followed by root-aligned
blending and joint-dynamics and collision constraints.

Strict leave-one-music-out evaluation achieves genre Recall@1 of 0.317 and
Recall@5 of 0.783, indicating stronger shortlist relevance than top-one
identification. In causal headless evaluation, all 135 formal runs, comprising
130 matched full-system and authored-timing short runs and five approximately
ten-minute full-system runs, satisfy the registered timing, readiness, and final
kinematic-output criteria at 60 Hz. The worst run-level control-work p99 is
14.483 ms, below the 16.67 ms budget, with no residual clearance or joint-range
violations. A separate 120 Hz pilot passes only six of 21 runs. Harmonic beat
alignment score relative to delivered online beat events is 0.313 with beat
synchronisation and 0.306 with authored timing; the difference is not
statistically significant after multiplicity correction.

The results establish an inspectable integration of music retrieval and
continuous humanoid playback with an empirically validated 60 Hz operating
point. Musical generalisation and the benefit of beat synchronisation remain
limited or unconfirmed. Validation concerns kinematic playback in MuJoCo;
dynamic balance, physical robot execution, and human-perceived dance quality
require further evaluation.
 before Chapter 1 in the parent thesis.
% Numerical claims follow the frozen Chapter 6 60 Hz v5 formal results.
% This fragment has no dependency on chapter-specific result macros.

\chapter*{Abstract}

Live music-driven humanoid dance requires musical relevance, responsive timing,
continuous motion transitions, and feasible robot configurations to be maintained
under causal audio access and limited computation. This dissertation presents
Humanoid Matcher, a retrieval-and-control system that adapts a finite library of
dance trajectories to streaming music on a 29-degree-of-freedom Unitree G1 model.

The offline pipeline combines paired AIST++ music and human motion, extracts
salient key poses from whole-body velocity minima, and retargets SMPL sequences
through General Motion Retargeting. Provenance checks and MuJoCo kinematic
preflight produce a catalogue of 411 admitted motions associated with 60 music
identities and 348 audio segments. Online, six-second rolling audio windows are
gain-normalised and encoded using Discogs EffNet embeddings, style activations,
and rhythm--timbre descriptors. Hierarchical retrieval combines musical similarity
with tempo, key-pose-density, and activity compatibility. Temporal evidence and
bounded-hold selection govern motion changes, while asynchronous retrieval and
candidate preparation support continuous playback. Motion-aware beat acceptance
and bounded phase correction adapt timing gradually. Transition-entry search
uses joint pose, velocity, foot contact, and root state, followed by root-aligned
blending and joint-dynamics and collision constraints.

Strict leave-one-music-out evaluation achieves genre Recall@1 of 0.317 and
Recall@5 of 0.783, indicating stronger shortlist relevance than top-one
identification. In causal headless evaluation, all 135 formal runs, comprising
130 matched full-system and authored-timing short runs and five approximately
ten-minute full-system runs, satisfy the registered timing, readiness, and final
kinematic-output criteria at 60 Hz. The worst run-level control-work p99 is
14.483 ms, below the 16.67 ms budget, with no residual clearance or joint-range
violations. A separate 120 Hz pilot passes only six of 21 runs. Harmonic beat
alignment score relative to delivered online beat events is 0.313 with beat
synchronisation and 0.306 with authored timing; the difference is not
statistically significant after multiplicity correction.

The results establish an inspectable integration of music retrieval and
continuous humanoid playback with an empirically validated 60 Hz operating
point. Musical generalisation and the benefit of beat synchronisation remain
limited or unconfirmed. Validation concerns kinematic playback in MuJoCo;
dynamic balance, physical robot execution, and human-perceived dance quality
require further evaluation.
 before Chapter 1 in the parent thesis.
% Numerical claims follow the frozen Chapter 6 60 Hz v5 formal results.
% This fragment has no dependency on chapter-specific result macros.

\chapter*{Abstract}

Live music-driven humanoid dance requires musical relevance, responsive timing,
continuous motion transitions, and feasible robot configurations to be maintained
under causal audio access and limited computation. This dissertation presents
Humanoid Matcher, a retrieval-and-control system that adapts a finite library of
dance trajectories to streaming music on a 29-degree-of-freedom Unitree G1 model.

The offline pipeline combines paired AIST++ music and human motion, extracts
salient key poses from whole-body velocity minima, and retargets SMPL sequences
through General Motion Retargeting. Provenance checks and MuJoCo kinematic
preflight produce a catalogue of 411 admitted motions associated with 60 music
identities and 348 audio segments. Online, six-second rolling audio windows are
gain-normalised and encoded using Discogs EffNet embeddings, style activations,
and rhythm--timbre descriptors. Hierarchical retrieval combines musical similarity
with tempo, key-pose-density, and activity compatibility. Temporal evidence and
bounded-hold selection govern motion changes, while asynchronous retrieval and
candidate preparation support continuous playback. Motion-aware beat acceptance
and bounded phase correction adapt timing gradually. Transition-entry search
uses joint pose, velocity, foot contact, and root state, followed by root-aligned
blending and joint-dynamics and collision constraints.

Strict leave-one-music-out evaluation achieves genre Recall@1 of 0.317 and
Recall@5 of 0.783, indicating stronger shortlist relevance than top-one
identification. In causal headless evaluation, all 135 formal runs, comprising
130 matched full-system and authored-timing short runs and five approximately
ten-minute full-system runs, satisfy the registered timing, readiness, and final
kinematic-output criteria at 60 Hz. The worst run-level control-work p99 is
14.483 ms, below the 16.67 ms budget, with no residual clearance or joint-range
violations. A separate 120 Hz pilot passes only six of 21 runs. Harmonic beat
alignment score relative to delivered online beat events is 0.313 with beat
synchronisation and 0.306 with authored timing; the difference is not
statistically significant after multiplicity correction.

The results establish an inspectable integration of music retrieval and
continuous humanoid playback with an empirically validated 60 Hz operating
point. Musical generalisation and the benefit of beat synchronisation remain
limited or unconfirmed. Validation concerns kinematic playback in MuJoCo;
dynamic balance, physical robot execution, and human-perceived dance quality
require further evaluation.
 before Chapter 1 in the parent thesis.
% Numerical claims follow the frozen Chapter 6 60 Hz v5 formal results.
% This fragment has no dependency on chapter-specific result macros.

\chapter*{Abstract}

Live music-driven humanoid dance requires musical relevance, responsive timing,
continuous motion transitions, and feasible robot configurations to be maintained
under causal audio access and limited computation. This dissertation presents
Humanoid Matcher, a retrieval-and-control system that adapts a finite library of
dance trajectories to streaming music on a 29-degree-of-freedom Unitree G1 model.

The offline pipeline combines paired AIST++ music and human motion, extracts
salient key poses from whole-body velocity minima, and retargets SMPL sequences
through General Motion Retargeting. Provenance checks and MuJoCo kinematic
preflight produce a catalogue of 411 admitted motions associated with 60 music
identities and 348 audio segments. Online, six-second rolling audio windows are
gain-normalised and encoded using Discogs EffNet embeddings, style activations,
and rhythm--timbre descriptors. Hierarchical retrieval combines musical similarity
with tempo, key-pose-density, and activity compatibility. Temporal evidence and
bounded-hold selection govern motion changes, while asynchronous retrieval and
candidate preparation support continuous playback. Motion-aware beat acceptance
and bounded phase correction adapt timing gradually. Transition-entry search
uses joint pose, velocity, foot contact, and root state, followed by root-aligned
blending and joint-dynamics and collision constraints.

Strict leave-one-music-out evaluation achieves genre Recall@1 of 0.317 and
Recall@5 of 0.783, indicating stronger shortlist relevance than top-one
identification. In causal headless evaluation, all 135 formal runs, comprising
130 matched full-system and authored-timing short runs and five approximately
ten-minute full-system runs, satisfy the registered timing, readiness, and final
kinematic-output criteria at 60 Hz. The worst run-level control-work p99 is
14.483 ms, below the 16.67 ms budget, with no residual clearance or joint-range
violations. A separate 120 Hz pilot passes only six of 21 runs. Harmonic beat
alignment score relative to delivered online beat events is 0.313 with beat
synchronisation and 0.306 with authored timing; the difference is not
statistically significant after multiplicity correction.

The results establish an inspectable integration of music retrieval and
continuous humanoid playback with an empirically validated 60 Hz operating
point. Musical generalisation and the benefit of beat synchronisation remain
limited or unconfirmed. Validation concerns kinematic playback in MuJoCo;
dynamic balance, physical robot execution, and human-perceived dance quality
require further evaluation.
